\section{A separate global optimizer: assigning invariant-preserving clocks}
\label{sec:sortingintervention}

This companion leaves the arbitrary finite-menu sampler architecture.
It optimizes all $n!$ assignments of kinetic resources to components using
local equilibrium and fixed-hold observations.  Its intervention mechanism
is more specific than a general repair or field change and must be kept
explicit.

\subsection{The declared operation and exact optimum}
There are $n$ independent repairable binary components.  Failure is the
single state $b=(1,\ldots,1)$.  Component $i$ has unknown equilibrium
failure probability $q_i\in(0,1)$.  The available resources are $n$ clock
modules with positive rates $r_1,\ldots,r_n$.  Installing module $j$ at
component $i$ produces the \emph{actual operating} rates
\begin{equation}
  0\longrightarrow1:\ r_jq_i,\qquad
  1\longrightarrow0:\ r_j(1-q_i).
  \label{eq:clockmodule}
\end{equation}
The module changes the refresh rate and preserves the equilibrium
probability.  This transferable factorization is structural physical
access, not an inference from a clock measured at an artificial bias.
It is neither a shared-repairman model nor a claim that every practical
repair enhancement has this form.  Every allowed assignment uses the same
multiset of modules and therefore the same declared resource inventory.

For a permutation $\sigma$, write $M_\sigma$ for the mean first failure
from $\pi(\cdot\mid b^c)$ and $p=\prod_iq_i$.  Put
\begin{align}
 S(q,r_\sigma)
 &=\int_0^\infty
 \left[\prod_{i=1}^n\{q_i+(1-q_i)e^{-r_{\sigma(i)}t}\}-p\right]dt\\
 &=\mathbb E\left[\frac1{\sum_i r_{\sigma(i)}B_i};
                        \ \sum_iB_i>0\right],
 \qquad B_i\sim\operatorname{Bernoulli}(1-q_i)
 \text{ independently}.
 \label{eq:weightedbern}
\end{align}
The second equality follows by expanding the product into its $2^n$
subsets and integrating their exponential factors; it is an identity, not
an instruction to enumerate those subsets.  The return identity gives
\[
 M_\sigma=\frac{S(q,r_\sigma)}{p(1-p)}.
\]

\paragraph{Theorem (exact assignment rule).}
Pair the $q_i$ and the $r_j$ in the \emph{same} sorted order.  This
assignment maximizes $M_\sigma$ over every permutation.  Opposite order
minimizes it.  Ties may be resolved arbitrarily.

\paragraph{Proof.}
Let $f_t(q,r)=q+(1-q)e^{-rt}$.  If $q_i\ge q_j$ and $a\ge b$, direct
algebra gives
\begin{equation}
 f_t(q_i,a)f_t(q_j,b)-f_t(q_i,b)f_t(q_j,a)
   =(q_i-q_j)(e^{-bt}-e^{-at})\ge0.
 \label{eq:sortswap}
\end{equation}
All other factors are nonnegative.  Removing inverted pairs therefore
increases the full return probability for \emph{every} $t\ge0$.
Subtract the common $p$, integrate, and divide by the common $p(1-p)$.
The reverse exchange gives the minimizing order.\hfill$\square$

This also proves an exact nuisance-elimination statement.  Adjoin any
independent finite ergodic spectator process, with any unknown invariant
law or timescales, and declare the target to be $(b,b_{\rm sp})$ for a
specified spectator state.  Its return probability multiplies both sides
of the exchange by a nonnegative function of time.  The same assignment
therefore remains optimal for the enlarged equilibrium-safe mean, without
learning the spectator.  This statement is exact for a certified ordering;
the uniform approximate-order bound below is for the product component
family and is not silently extended to arbitrarily slow spectators.
The spectator statement is conditional on the specified initial law.
Preparing that law can require its own slow relaxation, sensing, and
memory bill; no spectator preparation operation is supplied by the
ordering identity.

The endpoint is this particular mean objective and initial law.  The
exchange does not assert stochastic ordering of all first-failure times,
all terminal records, or survival from the all-working state.  Such stronger
claims do not follow merely from ordering return probabilities.

\subsection{Approximate order does not need $1/n$ probability precision}
Assume the supplied coarse bounds
\[
 0<q_i\le\bar q<1,\qquad
 0<r_*\le r_j\le r^*,\qquad \kappa=r^*/r_*.
\]
For $n\ge2$, define
\begin{equation}
 C_n(\bar q,\kappa)=
 \frac{3\kappa^2 n/[(n+1)(1-\bar q)^2]
          +\kappa n^2\bar q^{n-1}}
      {(1-\bar q^n)^2}.
 \label{eq:sortingconstant}
\end{equation}
This has a finite bound depending only on $\bar q$ and $\kappa$; for
example
\[
 C_n\le C_*:=
 \frac{3\kappa^2/(1-\bar q)^2
       +\kappa(1+\bar q)/(1-\bar q)^3}
      {(1-\bar q^2)^2}.
\]

\paragraph{Lemma (dimension-independent sensitivity).}
For any two probability vectors in $[0,\bar q]^n$ and fixed rate vector
in $[r_*,r^*]^n$,
\begin{equation}
 |\log S(q,r)-\log S(\widetilde q,r)|
       \le C_n\|q-\widetilde q\|_\infty.
 \label{eq:sortinglipschitz}
\end{equation}
Here $S$ remains positive even when an estimated $q_i$ is clipped to zero.
In addition, if $|\log(\widetilde r_i/r_i)|\le\rho$ for all $i$, then
$|\log S(q,\widetilde r)-\log S(q,r)|\le\rho$.

\paragraph{Proof.}
Set $F(0)=0$ and $F(w)=1/w$ for $w>0$, and let
$W_{-i}=\sum_{j\ne i}r_jB_j$.  Differentiating the finite product law gives
\[
 \partial_{q_i}S=\mathbb E[F(W_{-i})-F(W_{-i}+r_i)].
\]
For $W_{-i}>0$, the integrand equals
$r_i/[W_{-i}(W_{-i}+r_i)]$; for $W_{-i}=0$, it equals $-1/r_i$.
Writing $K=\sum_{j\ne i}B_j$ and $a=1-\bar q$, this gives
\[
 |\partial_{q_i}S|
 \le\frac{r^*}{r_*^2}
      \mathbb E\left[\frac1{K(K+1)};K>0\right]
       +\frac{\bar q^{n-1}}{r_*}.
\]
For all nonnegative integers $K$ the bracketed function is at most
$3/[(K+1)(K+2)]$.  This latter function is decreasing, and $K$
stochastically dominates $\operatorname{Bin}(n-1,a)$.  The beta integral
therefore yields
\begin{align*}
 \mathbb E\frac1{(K+1)(K+2)}
 &\le\int_0^1 t^{K_{\rm bin}}(1-t)\,dt\ \text{ averaged over }K_{\rm bin}\\
 &=\int_0^1(\bar q+at)^{n-1}(1-t)\,dt
 \le\frac1{a^2n(n+1)}.
\end{align*}
Consequently
\[
 \sum_i|\partial_{q_i}S|
 \le\frac{3r^*}{r_*^2(n+1)a^2}
           +\frac{n\bar q^{n-1}}{r_*}.
\]
For $W=\sum_i r_iB_i$, Cauchy--Schwarz gives
\[
 S\ge\frac{\Pr(W>0)^2}{\mathbb EW}
       \ge\frac{(1-\bar q^n)^2}{nr^*}.
\]
Dividing proves $\|\nabla\log S\|_1\le C_n$, hence
\eqref{eq:sortinglipschitz} by integrating along the segment.
The rate statement follows by coupling the same $B_i$ and comparing
each positive sum of rates multiplicatively.\hfill$\square$

\paragraph{Theorem (acquired near-optimal assignment).}
Suppose local evidence supplies simultaneous bounds
\[
 \|\widehat q-q\|_\infty\le\eta,\qquad
 |\log(\widehat r_j/r_j)|\le\rho,
\]
with $\widehat q$ clipped to $[0,\bar q]^n$.
Sort $\widehat q$ and $\widehat r$ in the same order and deploy the
resulting assignment $\widehat\sigma$.  Then
\begin{equation}
 \frac{M_{\widehat\sigma}}{\max_\sigma M_\sigma}
       \ge \exp[-2(C_n\eta+\rho)].
 \label{eq:acquiredsorting}
\end{equation}
Indeed the estimated sorting maximizes the estimated $S$ exactly.
Comparing true and estimated $S$ twice loses at most the displayed
factor.  The true denominator $p(1-p)$ is the same for every assignment
and is never estimated.  Choosing $\eta=\epsilon/(4C_n)$ and
$\rho=\epsilon/4$ gives an $e^{-\epsilon}$ guarantee.

\subsection{Finite local experiments at the actual operating context}
The probability and clock measurements can use fixed holds.  No physical
experiment must wait for the rare global target or measure a tiny first-exit
probability relatively.

For $q_i$, use any module with rate at least $r_*$.  From any prior bit,
a hold $t_q\ge r_*^{-1}\log(8/\eta)$ makes the conditional read mean
within $\eta/8$ of $q_i$.  With readout/hold/context errors contributing
another at most $\eta/8$, martingale Hoeffding on
\[
 m_q\ge\frac{8}{\eta^2}\log\frac{4n}{\delta}
\]
successive reads bounds the remaining error by $\eta/2$ simultaneously.
Clipping cannot increase the error.  Reusing a cell is allowed because the
bound is conditional on every previous experimental history; independent
raw readings are not assumed.  The source still must obey its declared
memoryless dynamics and faithful-context law.

To learn module $j$ without borrowing its clock from a biased experiment,
install it at a reference component in its actual operating environment.
Pin that component first to 0 and then, on a separate trial, to 1.  After
the same nominal hold $h=1/r^*$, the ideal success probabilities are
\[
 p_0=q(1-e^{-r_jh}),\qquad
 p_1=q+(1-q)e^{-r_jh},\qquad
 D_j=p_1-p_0=e^{-r_jh}.
\]
Thus $q$ cancels and $r_j=-h^{-1}\log D_j$.
For $0<\rho\le1/4$, let
$\alpha=\rho/(64e\kappa)$.
Estimate each of $p_0,p_1$ to absolute error at most $\alpha$, using
$O(\alpha^{-2}\log(n/\delta))$ trials per initial bit.  Include the
finite pin/read errors and the hold error in this allowance.  Since
$D_j\ge e^{-1}$, clipping the estimated difference to
$[e^{-1}/2,1]$ and finite interval evaluation of $-\log(D_j)/h$ give
absolute rate error less than $\rho r_*/4$, with ample slack for numerical
rounding.  Clipping the resulting rate to $[r_*,r^*]$ preserves the error.
This implies the desired log-rate error $\rho$.

For nonidentical actual hold durations, the probability error of either
read is at most $r^*\Delta_h$ when its timing deviation is at most
$\Delta_h$.  Taking $r^*\Delta_h\le\alpha/4$ is sufficient.
This is a fixed accuracy determined by the decision tolerance and local
clock contrast.  It does not shrink with the number of subsequent
simulated updates, because there are no such updates in this optimizer.
The two pins can be implemented serially.  An alternative simultaneous
two-cell exit race has total hazard $r_j$ only when independence of its
two physical clocks is actually supplied; that additional independence
is unnecessary here.

The source and deployment fields need not be numerically reconstructed.
They must reproduce the same declared operating environment, or have a
valid bound on the resulting discrepancy.  A common assignment-independent
probability drift $\|q^{\rm dep}-q^{\rm acq}\|_\infty\le\xi$ is included
by replacing $\eta$ by $\eta+\xi$ in
\eqref{eq:acquiredsorting}; it does not require $\xi=O(1/n)$.
An independently justified operating log-clock drift is similarly added
to $\rho$.  A bound on field digits alone does not justify a clock-drift
bound when the rate's field dependence is undeclared.  Assignment-dependent
changes to the equilibrium law invalidate the common-denominator
cancellation and require a different theorem.

At fixed $\bar q,\kappa$, the sufficient acquisition scales as
\begin{align*}
 \text{local reads and pins}\quad&
 O\!\left(n\epsilon^{-2}\log\frac n\delta\right),\\
 \text{serial free dwell}\quad&
 O\!\left(\frac n{r_*\epsilon^2}
     \log\frac1\epsilon\log\frac n\delta\right),\\
 \text{classical sorting}\quad& O(n\log n),
 \qquad \text{stored numbers}\quad O(n).
\end{align*}
The constants include $C_n^2$ and $\kappa^2$ and can be large.
Every actual context setup, device installation, priced pin, settle,
readout, timer, finite arithmetic, energy use, and module installation at
deployment is additional at its declared service cost.  This is a
finite theoretical operation; the formulas do not assert an economical
hardware implementation.

The statistical failure probability is additional to any failure
probability $\delta_{\rm cal}$ of a shared instrument or structural-margin
certificate.  On that common calibration event, the per-read conditional
bias bounds above are included directly in the estimator; they are not
charged by requiring every noisy read to be correct.  The joint acquired
assignment confidence is at least $1-\delta-\delta_{\rm cal}$.

An unresolved coarse probability or rate bound returns that status.
The transferable module law is supplied closure.  Short observations do
not establish that it holds at every untested site or exclude common
causes, shared repairs, ageing, or omitted failures.  An implementation
whose module changes $q_i$, or whose clock depends on the selected site
outside \eqref{eq:clockmodule}, is outside this optimizer.

\subsection{Preparation, separation, and prior comparison}
The acquired assignment concerns the equilibrium-safe target.
If additionally $q_i\ge q_*>0$ is certified and
$\min(q_*,1-\bar q)=q_{\min}$, then a product relaxation hold
\[
 T_{\rm prep}\ge r_*^{-1}\log\frac{4n}{\beta q_{\min}}
\]
gives a multiplicative equilibrium approximation uniformly over the
initial state.  Inspection and retaining only safe outcomes turn this
into the accepted density-ratio preparation, with the finite readout and
abort allowances charged.  Since $p\le\bar q^n$, a supplied safe-mass
floor $1-\bar q^n$ controls the retry bill.  If a preparation routine has
only TV accuracy, its error cannot simply be added to this unbounded-mean
objective.  The density ratio changes the deployed comparison by its
explicit multiplicative factor.  Module installation, safe preparation,
and the mission itself are physically separate operations.

The forecast lower bound in Section~\ref{sec:sameaccesslower} applies in
this same product family, even with perfect arbitrary-field experiments
and full local paths.  Thus near-optimal assignment among $n!$ choices can
use $\widetilde O(n)$ local evidence at fixed margins and accuracy while
an absolute mean forecast uniformly requires $\Omega(n^2)$ local
evidence.  This is a genuine information-task separation, not a comparison
with brute-force first-failure simulation.  The lower bound is polynomial,
not exponential; active local access already defeats inverse global
rarity for absolute forecasting.

Pairwise interchange and rearrangement methods are longstanding tools in
reliability allocation and stochastic scheduling.  The swap proof here
is an explicit application of that machinery.  We do not claim historical
priority for sorting allocations in reliability or for the Green-function
identity.  The earned result is the precise stationary-safe first-event
objective, its unknown-nuisance cancellation, the dimension-independent
approximation bound, and the complete actual-context acquisition and
same-access information comparison at the stated module scope.  Strong
ordinary methods using those same identities may realize the same
decision; a full simulator is not a necessary baseline.
